\documentclass[11pt,a4paper]{article}

% ── Codificación y fuentes ────────────────────────────────────────────────────
\usepackage[utf8]{inputenc}
\usepackage[T1]{fontenc}
\usepackage{lmodern}
\usepackage[spanish]{babel}

% ── Geometría y espaciado ─────────────────────────────────────────────────────
\usepackage[top=2.5cm, bottom=2.5cm, left=2.8cm, right=2.8cm]{geometry}
\usepackage{setspace}
\setstretch{1.15}
\usepackage{parskip}

% ── Tablas ────────────────────────────────────────────────────────────────────
\usepackage{booktabs}
\usepackage{tabularx}
\usepackage{array}
\usepackage{longtable}
\usepackage{enumitem}

% ── Colores y cajas ───────────────────────────────────────────────────────────
\usepackage[dvipsnames]{xcolor}
\usepackage{tcolorbox}
\tcbuselibrary{skins, breakable}

% ── Cabeceras y pies de página ────────────────────────────────────────────────
\usepackage{fancyhdr}
\usepackage{lastpage}
\pagestyle{fancy}
\fancyhf{}
\fancyhead[L]{\small\textbf{Informe de Evaluación} --- Física para Bioanalistas}
\fancyhead[R]{\small Barbara Marcano}
\fancyfoot[C]{\small Página \thepage\ de \pageref{LastPage}}
\renewcommand{\headrulewidth}{0.4pt}

% ── Hiperenlaces ──────────────────────────────────────────────────────────────
\usepackage[colorlinks=true, linkcolor=NavyBlue, urlcolor=NavyBlue]{hyperref}

% ── Paleta de colores personalizados ─────────────────────────────────────────
\definecolor{desempeno_excelente}{HTML}{1A6B2F}
\definecolor{desempeno_bueno}{HTML}{2E5A9C}
\definecolor{desempeno_suficiente}{HTML}{8B6914}
\definecolor{desempeno_deficiente}{HTML}{C0392B}
\definecolor{desempeno_insuficiente}{HTML}{7B241C}
\definecolor{grisFondo}{HTML}{F5F5F5}
\definecolor{azulTitulo}{HTML}{1B2A6B}
\definecolor{verdeFortaleza}{HTML}{1A6B2F}
\definecolor{naranjaMejora}{HTML}{A04000}

% ── Estilos de cajas ─────────────────────────────────────────────────────────
\tcbset{
  cajaTitulo/.style={
    enhanced, breakable,
    colback=azulTitulo!8, colframe=azulTitulo,
    fonttitle=\bfseries\large, coltitle=white,
    attach boxed title to top left={yshift=-2mm, xshift=4mm},
    boxed title style={colback=azulTitulo},
    top=6pt, bottom=6pt, left=8pt, right=8pt,
  },
  cajaFortaleza/.style={
    enhanced, breakable,
    colback=verdeFortaleza!6, colframe=verdeFortaleza!60,
    fonttitle=\bfseries, coltitle=verdeFortaleza,
    top=4pt, bottom=4pt, left=8pt, right=8pt,
  },
  cajaMejora/.style={
    enhanced, breakable,
    colback=naranjaMejora!6, colframe=naranjaMejora!60,
    fonttitle=\bfseries, coltitle=naranjaMejora,
    top=4pt, bottom=4pt, left=8pt, right=8pt,
  },
  cajaRecomendacion/.style={
    enhanced, breakable,
    colback=grisFondo, colframe=gray!50,
    fonttitle=\bfseries, coltitle=black,
    top=4pt, bottom=4pt, left=8pt, right=8pt,
  },
}

% ─────────────────────────────────────────────────────────────────────────────
\begin{document}

% ══ PORTADA ══════════════════════════════════════════════════════════════════
\begin{center}
  {\Large\bfseries\color{azulTitulo} Universidad de Oriente/Núcleo Sucre --- Física para Bioanalistas}\\[4pt]
  {\large Informe de Evaluación Preliminar}\\[12pt]
  \rule{\linewidth}{1.2pt}\\[6pt]
  {\LARGE\bfseries Barbara Marcano}\\[4pt]
  \rule{\linewidth}{0.6pt}
\end{center}

% ══ FICHA TÉCNICA ════════════════════════════════════════════════════════════
\vspace{4pt}
\begin{tcolorbox}[cajaTitulo, title=Datos del informe]
\begin{tabular}{@{}llll@{}}
  \textbf{Estudiante:}  & Barbara Marcano  &
  \textbf{Fecha:}       & 2026-08-08 \\[4pt]
  \textbf{Archivo PDF:} & \texttt{Barbara\_Marcano.pdf}  &
  \textbf{Páginas:}     & 3 \\
\end{tabular}
\end{tcolorbox}

% ══ RESULTADO GLOBAL ═════════════════════════════════════════════════════════
\vspace{6pt}
\begin{center}
  \begin{tcolorbox}[
    enhanced,
    colback=white, colframe=desempeno_deficiente,
    width=0.62\linewidth,
    boxrule=2pt,
    halign=center, valign=center,
    top=8pt, bottom=8pt,
  ]
    \centering
    {\large\bfseries Puntaje sugerido}\\[6pt]
    {\Huge\bfseries\color{desempeno_deficiente} 5.9/10}\\[6pt]
    {\large Nivel de desempeño: \textbf{\color{desempeno_deficiente} Deficiente}}
  \end{tcolorbox}
\end{center}

% ══ RESUMEN GENERAL ══════════════════════════════════════════════════════════
\section*{Resumen general}
El estudiante demuestra una comprensión parcial de los conceptos evaluados en la unidad de Líquidos y Fenómenos de Transporte. Si bien acierta en algunas explicaciones cualitativas y resultados numéricos, presenta deficiencias significativas en la justificación de sus respuestas, la demostración de procedimientos de cálculo y, en particular, en la comprensión de principios fundamentales como la ósmosis y la capilaridad. La falta de argumentación detallada y la omisión de pasos en los cálculos son recurrentes, lo que impacta negativamente en la evaluación.

% ══ FORTALEZAS Y ÁREAS DE MEJORA ════════════════════════════════════════════
\vspace{4pt}
\begin{minipage}[t]{0.48\linewidth}
  \begin{tcolorbox}[cajaFortaleza, title={Fortalezas}]
\begin{itemize}[leftmargin=1.5em, itemsep=2pt]
  \item Correcta identificación del efecto de la humedad en la evaporación y disipación de calor (Pregunta 1b).
  \item Buena comprensión del mecanismo de colapso alveolar y el papel del surfactante (Pregunta 2a y 2b).
  \item Entendimiento del impacto del espesor de la membrana en la eficiencia de difusión y su integridad estructural en hemodiálisis (Pregunta 3a y 3b).
  \item Correcta identificación del efecto de una solución hipertónica en los eritrocitos (Pregunta 4b).
  \item Cálculos numéricos correctos para la disipación de calor (Pregunta 1c), la tasa de difusión (Pregunta 3c) y la presión osmótica (Pregunta 4c), aunque sin mostrar el procedimiento.
\end{itemize}
  \end{tcolorbox}
\end{minipage}
\hfill
\begin{minipage}[t]{0.48\linewidth}
  \begin{tcolorbox}[cajaMejora, title={Areas de mejora}]
\begin{itemize}[leftmargin=1.5em, itemsep=2pt]
  \item Necesidad de argumentar y justificar completamente las respuestas cualitativas, explicando los principios físicos subyacentes.
  \item Obligatoriedad de mostrar todos los pasos de los cálculos, incluyendo fórmulas, sustitución de valores y unidades.
  \item Revisión profunda de conceptos fundamentales de ósmosis y el movimiento del agua a través de membranas semipermeables.
  \item Diferenciación clara entre distintos fenómenos de la mecánica de fluidos, como la capilaridad y el flujo de fluidos por tuberías.
  \item Atención a la lectura de las preguntas para aplicar la ley física correcta al escenario planteado.
\end{itemize}
  \end{tcolorbox}
\end{minipage}

% ══ OBSERVACIONES POR PREGUNTA ═══════════════════════════════════════════════
\section*{Observaciones por pregunta}

\begin{longtable}{>{\bfseries}p{2.8cm} p{1.6cm} p{7.0cm} p{2.8cm}}
  \toprule
  \textbf{Pregunta} & \textbf{Puntaje} & \textbf{Evaluación} & \textbf{Errores detectados} \\
  \midrule
  \endhead
  \bottomrule
  \endfoot
    \textbf{Pregunta 1: Termorregulación y Eficiencia de la Sudoración} & 1.8/2.0 & a) La respuesta es correcta en su conclusión (paciente B necesita más sudor), pero carece de la justificación física completa sobre por qué el sudor que gotea es menos eficiente que el que se evapora (calor latente de vaporización vs. calor sensible). b) La explicación es excelente, detallando el efecto de la humedad en la evaporación y el riesgo de golpe de calor. c) El cálculo es correcto y las unidades son consistentes. & \textit{Explicación incompleta del mecanismo de eficiencia de la sudoración en 1a.} \\
    \midrule
    \textbf{Pregunta 2: Mecánica Respiratoria y Ley de Laplace} & 1.4/2.0 & a) La explicación es correcta, aplicando la Ley de Laplace para explicar el colapso de alvéolos pequeños. b) La explicación sobre el papel del surfactante es correcta y clara. c) Se proporciona solo el resultado numérico, el cual es incorrecto. No se muestra ningún procedimiento de cálculo, lo cual es una falta grave según los criterios de evaluación. & \textit{Resultado numérico incorrecto en 2c; Ausencia de procedimiento de cálculo en 2c.} \\
    \midrule
    \textbf{Pregunta 3: Optimización en Hemodiálisis y Primera Ley de Fick} & 1.7/2.0 & a) La aplicación de la Primera Ley de Fick para determinar el aumento de eficiencia es correcta, indicando un aumento del 40\%. b) La justificación sobre los riesgos de una membrana demasiado delgada es correcta y pertinente. c) Se proporciona solo el resultado numérico, el cual es correcto. Sin embargo, no se muestra ningún procedimiento de cálculo, lo cual penaliza el puntaje. & \textit{Ausencia de procedimiento de cálculo en 3c.} \\
    \midrule
    \textbf{Pregunta 4: Errores de Infusión Intravenosa y Ósmosis} & 1.0/2.0 & a) La respuesta es conceptualmente incorrecta. El estudiante identifica correctamente el agua destilada como hipotónica, pero afirma erróneamente que el agua 'sale' de los eritrocitos, cuando en realidad 'entra', causando hemólisis. Este es un error conceptual fundamental. b) La explicación sobre el efecto de una solución hipertónica (NaCl 5\%) es correcta, incluyendo la crenación. c) Se proporciona solo el resultado numérico, el cual es correcto. Sin embargo, no se muestra ningún procedimiento de cálculo, lo cual penaliza el puntaje. & \textit{Error conceptual fundamental en 4a (dirección del flujo de agua en solución hipotónica).; Ausencia de procedimiento de cálculo en 4c.} \\
    \midrule
    \textbf{Pregunta 5: Capilaridad y Toma de Muestras Sanguíneas} & 0.0/2.0 & a) y b) Las respuestas están en blanco. c) El estudiante confunde el problema de capilaridad con un problema de flujo de fluidos (Ley de Poiseuille), aplicando una fórmula incorrecta y calculando una magnitud diferente (caudal 'Q' en lugar de altura 'h'). Además, los cálculos presentados para 'Q' contienen errores e inconsistencias. La pregunta está completamente errada en su enfoque. & \textit{Respuestas en blanco en 5a y 5b.; Error conceptual grave en 5c (aplicación de ley física incorrecta).; Errores procedimentales en el cálculo de 5c.} \\
    \midrule
\end{longtable}

% ══ ERRORES DETECTADOS ═══════════════════════════════════════════════════════
\section*{Análisis de errores}

\subsection*{Errores conceptuales}
\begin{itemize}[leftmargin=1.5em, itemsep=2pt]
  \item Malentendido fundamental sobre la dirección del flujo de agua en ósmosis cuando una célula se expone a una solución hipotónica (Pregunta 4a).
  \item Confusión entre el fenómeno de capilaridad y el flujo de fluidos por tuberías (Ley de Poiseuille) en la Pregunta 5c, lo que llevó a aplicar una ley física completamente incorrecta.
  \item Explicación incompleta del mecanismo de eficiencia de la sudoración en la Pregunta 1a, al no diferenciar claramente entre la disipación de calor por evaporación y por goteo.
\end{itemize}

\subsection*{Errores procedimentales}
\begin{itemize}[leftmargin=1.5em, itemsep=2pt]
  \item Omisión sistemática de los pasos de cálculo en las Preguntas 2c, 3c y 4c, presentando solo el resultado final.
  \item Resultado numérico incorrecto en la Pregunta 2c.
  \item Cálculos inconsistentes y erróneos en la Pregunta 5c, además de aplicar la fórmula incorrecta.
\end{itemize}

% ══ RECOMENDACIONES AL ESTUDIANTE ════════════════════════════════════════════
\section*{Retroalimentación para el estudiante}
\begin{tcolorbox}[cajaRecomendacion, title={Dirigido a: Barbara Marcano}]
Estimado/a estudiante, es crucial que profundices en la comprensión de los principios físicos subyacentes a cada fenómeno. Para las preguntas cualitativas, no basta con dar una respuesta; debes explicar el 'por qué' y 'cómo' ocurre, justificando tu razonamiento con los conceptos aprendidos. En los problemas cuantitativos, es indispensable que muestres todos los pasos de tus cálculos: la fórmula utilizada, la sustitución de los valores con sus unidades y el resultado final con las unidades correctas. Presta especial atención a la lectura de cada problema para identificar correctamente el fenómeno físico que se te pide analizar. Te recomiendo revisar a fondo los temas de ósmosis y capilaridad, ya que se identificaron errores conceptuales importantes en estas áreas.
\end{tcolorbox}

% ══ NOTA PARA EL PROFESOR ════════════════════════════════════════════════════
\section*{Nota para el profesor}
\begin{tcolorbox}[
  enhanced, breakable,
  colback=yellow!8, colframe=yellow!60!black,
  fonttitle=\bfseries, coltitle=black,
  title={Observaciones para revisión manual},
  top=4pt, bottom=4pt, left=8pt, right=8pt,
]
El estudiante presenta un desempeño deficiente, con una puntuación de 5.9/10.0. Si bien hay aciertos puntuales, la falta de justificación y la omisión de procedimientos de cálculo son generalizadas, lo cual va en contra del criterio fundamental de evaluación. Destacan errores conceptuales graves en la Pregunta 4a (ósmosis) y la Pregunta 5 (confusión de capilaridad con flujo de fluidos). Sugiero enfatizar la importancia de la argumentación y la demostración de los pasos en los cálculos en futuras evaluaciones.
\end{tcolorbox}

% ══ PIE DE INFORME ════════════════════════════════════════════════════════════
\vspace{12pt}
\begin{center}
  \footnotesize\color{gray}
  Informe generado automáticamente por el sistema de evaluación con IA (gemini-2.5-flash) y soporte LaTex. \\
  Este documento es una evaluación \textbf{preliminar} para ser revisado por el profesor antes de ser entregado al 
  estudiante. \\
  Generado el 2026-08-08 \textbullet\ BIOANALISIS Sistema Académico de Evaluación.
  \end{center}

\end{document}
